\documentclass[11pt,a4paper]{article}
\usepackage[margin=25mm]{geometry}
\usepackage{iftex}
\ifPDFTeX
  \usepackage[T2A]{fontenc}
  \usepackage[utf8]{inputenc}
\else
  \usepackage{fontspec}
  \setmainfont{cmunrm.otf}[
    BoldFont=cmunbx.otf,
    ItalicFont=cmunti.otf,
    BoldItalicFont=cmunbi.otf,
    Ligatures=TeX]
\fi
\usepackage[english,russian]{babel}
\usepackage{amsmath,amssymb,amsthm}
\usepackage{microtype}
\usepackage{enumitem}
\usepackage[unicode]{hyperref}
\hypersetup{
  pdftitle={Расширения квандлов со стабилизирующимися орбитальными рядами},
  pdfauthor={Danil Pokulevskii},
  pdflang={ru-RU},
  colorlinks=true,
  linkcolor=blue,
  citecolor=blue,
  urlcolor=blue}
\setlength{\emergencystretch}{2em}
\setlist[enumerate]{itemsep=2pt,topsep=5pt}
\frenchspacing

\newtheorem{theorem}{Теорема}[section]
\newtheorem{lemma}[theorem]{Лемма}
\newtheorem{proposition}[theorem]{Предложение}
\newtheorem{corollary}[theorem]{Следствие}
\theoremstyle{definition}
\newtheorem{definition}[theorem]{Определение}
\theoremstyle{remark}
\newtheorem{remark}[theorem]{Замечание}
\renewcommand{\proofname}{Доказательство}
\newcommand{\OS}{\mathcal{OS}}
\newcommand{\tOS}{t\mathcal{OS}}
\newcommand{\Inn}{\operatorname{Inn}}
\newcommand{\Orb}{\operatorname{Orb}}
\newcommand{\Aut}{\operatorname{Aut}}
\newcommand{\Z}{\mathbb Z}
\newcommand{\conj}{\triangleright}

\title{Расширения квандлов\\со стабилизирующимися орбитальными рядами}
\author{Danil Pokulevskii\\[3pt]
  \small Новосибирский государственный технический университет\\
  \small Новосибирск, Россия\\
  \small\href{mailto:555tery@gmail.com}{\texttt{555tery@gmail.com}}}
\date{7 октября 2026 года\\[2pt]\small Рабочая версия препринта}

\begin{document}
\maketitle

\begin{abstract}
Пусть класс \(\OS\) состоит из квандлов, у которых каждый орбитальный
ряд стабилизируется после конечного числа шагов. Доказывается, что
\(\OS\) замкнут относительно расширений: если фактор и каждый полный
слой эпиморфизма принадлежат \(\OS\), то этому классу принадлежит и
исходный квандл. Никакой общей границы глубин слоёв не требуется.
Ключевым является групповое тождество для относительных коммутаторных
рядов одного автоморфизма. Если \(N\trianglelefteq L\),
\(\phi(N)=N\) и \([L/N,\bar\phi]=L/N\), то ряды
\(R_0=N\), \(R_{i+1}=[R_i,\phi]\) и
\(D_0=L\), \(D_{i+1}=[D_i,\phi]\) удовлетворяют
\(D_i=R_iD_{i+1}\). Переход к квандлам проводится через классы
сопряжения и накрытия и не предполагает инъективности отображения
в левую трансляцию. Отдельно показывается, что класс
\(\OS_\omega=\bigcup_{n\ge1}\OS_n\) не замкнут относительно расширений.
Тем самым получены ответы на обе части вопроса~5.7
Бонатто--Кранс--Насыбуллова--Уитни.
\end{abstract}

\noindent\textbf{Ключевые слова:} квандл, орбитальный ряд, расширение,
субнормальная подгруппа, автоморфизм группы.

\section{Введение и основные результаты}

Орбитальные ряды квандлов и классы \(\OS\), \(\OS_n\) введены в
\cite{BCNW}. Вопрос~5.7 этой работы спрашивает о замкнутости
\(\OS\) и \(\OS_\omega\) относительно расширений.
Нумерация вопроса и результатов \cite{BCNW}, используемая ниже,
относится к версии arXiv v3 от 18 августа 2020 года.

Принципиально важно различать два условия. В классе \(\OS\)
каждая отдельная орбитальная цепь должна стабилизироваться, но число
шагов может зависеть от цепи. В классе \(\OS_\omega\) одна конечная
граница должна работать для всех цепей данного квандла.
В определении расширения общая граница по разным полным слоям
не предполагается.

\begin{theorem}\label{thm:main}
Пусть \(\pi:Q\twoheadrightarrow B\) --- эпиморфизм квандлов.
Если \(B\in\OS\) и \(\pi^{-1}(b)\in\OS\) для каждого \(b\in B\),
то \(Q\in\OS\). Следовательно, класс \(\OS\) замкнут относительно
расширений.
\end{theorem}

\begin{proposition}\label{prop:negative}
Класс \(\OS_\omega\) не замкнут относительно расширений.
Контрпример можно выбрать с тривиальным основанием и конечными
полными слоями.
\end{proposition}

Теорема~\ref{thm:main} доказывается в разделе~\ref{sec:quandles}.
Её групповое ядро, теорема~\ref{thm:group}, утверждает, что
стабилизация относительного коммутаторного ряда на нормальной
подгруппе поднимается на всю группу без увеличения числа шагов,
если фактор относительно совершенен для рассматриваемого
автоморфизма. Доказательство использует элементарный факт:
субнормальная подгруппа конечного дефекта, нормально порождающая
группу, совпадает с ней.
Предложение~\ref{prop:negative} следует из дизъюнктной суммы
диэдральных квандлов из примера~4.2 работы \cite{BCNW}
после рассмотрения проекции на номер слагаемого.

\section{Определения и вспомогательные леммы}\label{sec:prelim}

\begin{definition}
\emph{Квандл} --- непустое множество \(Q\) с операцией \(\conj\),
для которой
\[
 x\conj x=x,\qquad
 L_x:y\longmapsto x\conj y\text{ биективно},\qquad
 x\conj(y\conj z)=(x\conj y)\conj(x\conj z).
\]
Под-квандл замкнут относительно операции и обратных левых трансляций.
Группа \(\Inn(Q)=\langle L_x:x\in Q\rangle\) называется
внутренней группой квандла. Через \(\Orb(x,Q)\) обозначается
её орбита точки \(x\). Квандл связен, если эта группа транзитивна;
квандл тривиален, если \(x\conj y=y\) для всех \(x,y\).
\end{definition}

Выбирая последовательно \(x_i\in Q_i\), строим \emph{орбитальный ряд}
\[
 Q_0=Q,\qquad Q_{i+1}=\Orb(x_i,Q_i).
\]
Класс \(\OS\) состоит из квандлов, для каждого орбитального ряда
которых существует \(m\ge0\) с \(Q_m=Q_{m+1}\).
Класс \(\OS_n\), \(n\ge1\), задаётся требованием выбрать такое
\(m\le n\) для каждого ряда. Полагаем
\(\OS_\omega=\bigcup_{n\ge1}\OS_n\).
Стабильный член ряда связен, поскольку
\(Q_m=\Orb(x_m,Q_m)\).
Тривиализирующие классы \(\tOS\) и \(\tOS_n\) определяются
требованием достижения одноэлементного квандла, соответственно
за конечное число шагов или не позднее шага \(n\).

Под \emph{расширением} понимаем эпиморфизм
\(\pi:Q\twoheadrightarrow B\); его полные слои суть
\(F_b=\pi^{-1}(b)\). Это равносильно рассмотрению фактора
\(Q/\alpha\) и всех полных классов конгруэнции \(\alpha=\ker\pi\),
как в \cite[раздел~2.2]{BCNW}.

\begin{lemma}\label{lem:images}
Эпиморфизм \(p:E\twoheadrightarrow C\) индуцирует эпиморфизм
\(\Inn(E)\twoheadrightarrow\Inn(C)\), переводящий \(L_x\)
в \(L_{p(x)}\), и
\[
 p(\Orb(x,E))=\Orb(p(x),C).
\]
В частности, образы орбитального ряда образуют орбитальный ряд,
а класс \(\OS\) замкнут относительно эпиморфных образов.
\end{lemma}
\begin{proof}
Равенство \(pL_x=L_{p(x)}p\) верно также для обратных трансляций.
Если слово из трансляций тождественно на \(E\), то соответствующее
слово тождественно на \(C\), поскольку \(p\) сюръективно.
Это даёт корректность группового отображения, а сюръективность
следует из подъёма генераторов. Отсюда получается равенство орбит.
К произвольному под-квандлу применяем тот же аргумент для ограничения
на его образ.

Любой орбитальный ряд в \(C\) последовательно поднимается до ряда
в \(E\): очередной якорь поднимается в текущий член, который
сюръективно отображается на соответствующий член в \(C\).
Стабилизация поднятого ряда даёт стабилизацию исходного.
\end{proof}

\begin{lemma}\label{lem:invariant}
Если \(F\in\OS\), а непустой под-квандл \(S\subseteq F\)
инвариантен относительно \(\Inn(F)\), то \(S\in\OS\).
\end{lemma}
\begin{proof}
Возьмём любой ряд \(S=S_0\supseteq S_1\supseteq\cdots\)
с якорями \(x_j\in S_j\). По тем же якорям определим
\(F_0=F\), \(F_{j+1}=\Orb(x_j,F_j)\).
Индукцией получаем
\[
 F_j\supseteq S_j\supseteq F_{j+1},
 \qquad S_j\text{ инвариантен относительно }\Inn(F_j).
\]
Действительно, из инвариантности \(S_j\) следует
\(F_{j+1}\subseteq S_j\), а из \(S_j\subseteq F_j\) ---
\(S_{j+1}\subseteq F_{j+1}\).
Поскольку \(S_{j+1}\) есть орбита в \(S_j\), а
\(F_{j+1}\subseteq S_j\), трансляции от элементов \(F_{j+1}\)
и их обратные сохраняют \(S_{j+1}\). Это даёт следующий шаг.

Ряд \(F_j\) стабилизируется:
\(F_m=F_{m+1}=\Orb(x_m,F_m)\).
По включениям \(S_m=F_m\); связность \(F_m\) даёт
\(S_{m+1}=S_m\).
\end{proof}

\begin{remark}
В лемме~\ref{lem:invariant} существенно требование инвариантности.
Наследование \(\OS\) произвольными под-квандлами не используется;
в общем случае оно неверно \cite[пример~5.3]{BCNW}.
\end{remark}

\begin{lemma}\label{lem:slices}
Пусть все полные слои эпиморфизма
\(\pi:Q\twoheadrightarrow B\) принадлежат \(\OS\),
а \(Q_i\) --- орбитальный ряд в \(Q\).
Тогда каждый непустой срез \(Q_i\cap F_b\) принадлежит \(\OS\).
\end{lemma}
\begin{proof}
Если \(S_i=Q_i\cap F_b\ne\varnothing\), то непусты все
\(S_j=Q_j\cap F_b\), \(0\le j\le i\).
Для \(j<i\) и \(y\in S_j\) трансляции \(L_y^{\pm1}\)
сохраняют орбиту \(Q_{j+1}\) в \(Q_j\) и слой \(F_b\):
на основании \(L_b^{\pm1}\) фиксируют \(b\).
Поэтому \(S_{j+1}\) инвариантен относительно \(\Inn(S_j)\).
Начав с \(S_0=F_b\), последовательно применяем
лемму~\ref{lem:invariant} до индекса \(i\).
Непустота последующих срезов не требуется.
\end{proof}

\begin{lemma}\label{lem:trivialbase}
Эпиморфизм на тривиальный квандл с полными слоями из \(\OS\)
имеет источник в \(\OS\).
\end{lemma}
\begin{proof}
Первая орбита любого ряда лежит в одном полном слое.
Как орбита внутренней группы источника она инвариантна относительно
внутренней группы этого слоя. Применяем
лемму~\ref{lem:invariant}; оставшийся ряд стабилизируется.
\end{proof}

\begin{definition}
Эпиморфизм \(p:E\twoheadrightarrow C\) называется
\emph{накрытием}, если
\[
 p(x)=p(y)\quad\Longrightarrow\quad L_x^E=L_y^E.
\]
\end{definition}

\begin{lemma}\label{lem:covering}
Если \(p:E\twoheadrightarrow C\) --- накрытие и \(C\in\OS\),
то \(E\in\OS\).
\end{lemma}
\begin{proof}
Для ряда \(E_i\) образы \(C_i=p(E_i)\) стабилизируются
на связном под-квандле \(C_*\). На соответствующем хвосте все
\(E_i\twoheadrightarrow C_*\).
Рассмотрим группы глобальных перестановок множества \(E\)
\[
 \Gamma_i=\langle L_x^E:x\in E_i\rangle.
\]
Множества их генераторов одинаковы на хвосте: трансляция
зависит только от \(p(x)\), а множества образов \(p(E_i)=C_*\)
совпадают. Обозначим общую группу через \(\Gamma\).
Каждый \(E_i\) сохраняется этой группой, а её ограничения
порождают \(\Inn(E_i)\).
Поэтому \(E_{i+1}\) есть \(\Gamma\)-орбита;
следующий якорь принадлежит этой же орбите, и
\(E_{i+2}=E_{i+1}\).
\end{proof}

\section{Подъём относительного коммутаторного ряда}\label{sec:groups}

Для группы пишем \(x^y=y^{-1}xy\).
Если \(\phi\in\Aut(L)\), а \(X\le L\) удовлетворяет
\(\phi(X)=X\), положим
\[
 [X,\phi]=\langle x^{-1}\phi(x):x\in X\rangle.
\]
Здесь порождается подгруппа, а не только множество значений.

\begin{lemma}\label{lem:subnormal}
Пусть \(A\) субнормальна конечного дефекта в \(B\).
Если \(\langle A^B\rangle=B\), то \(A=B\).
\end{lemma}
\begin{proof}
Выберем конечную цепь
\(A=A_0\trianglelefteq A_1\trianglelefteq\cdots
\trianglelefteq A_m=B\).
Если \(m=0\), всё доказано.
Иначе нормальная подгруппа \(A_{m-1}\) группы \(B\) содержит \(A\),
а значит, содержит её нормальное замыкание \(B\).
Поэтому \(A_{m-1}=B\). Снимая последние звенья, получаем \(A=B\).
\end{proof}

\begin{theorem}\label{thm:group}
Пусть \(N\trianglelefteq L\), \(\phi\in\Aut(L)\),
\(\phi(N)=N\), а индуцированный автоморфизм \(\bar\phi\)
удовлетворяет \([L/N,\bar\phi]=L/N\).
Определим
\[
 R_0=N,\quad R_{i+1}=[R_i,\phi],\qquad
 D_0=L,\quad D_{i+1}=[D_i,\phi].
\]
Тогда при каждом \(i\ge0\)
\begin{equation}\label{eq:strong}
 D_i=R_iD_{i+1}.
\end{equation}
В частности,
\[
 R_r=R_{r+1}\quad\Longrightarrow\quad D_r=D_{r+1}
 \qquad(r\ge0).
\]
\end{theorem}
\begin{proof}
Обозначим \(\delta(x)=x^{-1}\phi(x)\). Имеем
\begin{equation}\label{eq:delta}
 \delta(xy)=\delta(x)^y\delta(y),\qquad
 \delta(x)^y=\delta(xy)\delta(y)^{-1}.
\end{equation}
Отсюда \([X,\phi]\trianglelefteq X\) для \(\phi\)-инвариантной \(X\).
Равенство \(\phi(\delta(x))=\delta(\phi(x))\) показывает
\(\phi\)-инвариантность этой подгруппы.
Ряды \(R_i,D_i\) убывают; монотонность операции
\(X\mapsto[X,\phi]\) даёт \(R_i\le D_i\).
Кроме того, \(R_i\) субнормальна в \(D_i\):
пересекаем конечную цепь
\[
 R_i\trianglelefteq R_{i-1}\trianglelefteq\cdots
 \trianglelefteq N\trianglelefteq L
\]
с \(D_i\). Нормальность соседних звеньев сохраняется;
крайние члены полученной цепи суть \(R_i\) и \(D_i\).

Докажем \eqref{eq:strong} индукцией.
При \(i=0\) образ \(D_1\) в \(L/N\) равен всему фактору,
поэтому \(L=ND_1\).
Пусть \(i\ge1\) и \(D_{i-1}=R_{i-1}D_i\).
Положим \(M_i=\langle R_i^{D_i}\rangle\).
Любой \(x\in D_{i-1}\) имеет вид \(rd\),
где \(r\in R_{i-1}\), \(d\in D_i\).
По \eqref{eq:delta}
\[
 \delta(x)=\delta(r)^d\delta(d)\in M_iD_{i+1}.
\]
Так как \(D_{i+1}\trianglelefteq D_i\), справа стоит подгруппа.
Все генераторы \(D_i=[D_{i-1},\phi]\) лежат в ней,
а обратное включение очевидно. Получаем
\[
 D_i=M_iD_{i+1}.
\]
Образ \(R_i\) в \(D_i/D_{i+1}\) субнормален конечного дефекта
и нормально порождает этот фактор.
По лемме~\ref{lem:subnormal} он равен всему фактору,
что даёт \eqref{eq:strong} и завершает индукцию.

Если \(R_r=R_{r+1}\), то
\(R_r=[R_r,\phi]\le[D_r,\phi]=D_{r+1}\).
Теперь \eqref{eq:strong} даёт \(D_r=D_{r+1}\).
\end{proof}

\begin{remark}
В теореме~\ref{thm:group} не предполагаются конечность,
разрешимость, расщепление или конечный порядок автоморфизма.
Члены \(R_i\) не обязаны быть нормальными во всей \(L\);
используется только их субнормальность в \(D_i\).
\end{remark}

\section{Расширения классов сопряжения}\label{sec:conjugacy}

Класс \(c^K=\{g^{-1}cg:g\in K\}\) рассматриваем как квандл
с операцией \(a\conj b=aba^{-1}\).
Это соглашение соответствует \(\mathrm{Conj}_{-1}\) в \cite{BCNW}.

\begin{lemma}\label{lem:towers}
Для \(c\in K\) положим \(K_0=K\),
\(K_{j+1}=\langle c^{K_j}\rangle\).
Квандл \(c^K\) принадлежит \(\OS\) тогда и только тогда,
когда ряд подгрупп \(K_j\) стабилизируется за конечное число шагов.
\end{lemma}
\begin{proof}
Положим \(Y_j=c^{K_j}\).
Тогда \(K_{j+1}\trianglelefteq K_j\) и
\(Y_{j+1}=\Orb(c,Y_j)\), так что \(Y_j\) --- главный
орбитальный ряд с постоянным якорем \(c\).
Равенство \(K_j=K_{j+1}\) даёт \(Y_j=Y_{j+1}\), а равенство
\(Y_j=Y_{j+1}\) даёт \(K_{j+1}=K_{j+2}\).

Покажем, что стабилизация главного ряда достаточна.
Любой конечный начальный отрезок произвольного орбитального
ряда можно сопряжением привести к \(Y_j\).
Если первые члены уже приведены, очередной якорь имеет вид
\(c^g\), где \(g\in K_j\).
Сопряжение на \(g^{-1}\) переводит его в \(c\)
и сохраняет все предыдущие \(Y_k\), \(k\le j\),
поскольку \(K_j\le K_k\).
Оно переводит следующую орбиту в \(Y_{j+1}\).
Таким образом, если \(Y_m=Y_{m+1}\), нормализация вплоть
до якоря с индексом \(m\) даёт равенство соответствующих
членов любого ряда. Обратное направление следует из
стабилизации главного ряда.
\end{proof}

\begin{lemma}\label{lem:kernelfiber}
Пусть \(\sigma:G\twoheadrightarrow H\), \(h=\sigma(c)\),
\(N=\ker\sigma\), \(P=\sigma^{-1}(C_H(h))\).
Если полный слой \(c^P\) эпиморфизма \(c^G\to h^H\)
принадлежит \(\OS\), то \(c^N\in\OS\).
\end{lemma}
\begin{proof}
Положим \(V=N\langle c\rangle=\sigma^{-1}(\langle h\rangle)\).
Так как \(\langle h\rangle\) центральна в \(C_H(h)\),
имеем \(V\trianglelefteq P\).
Каждый элемент \(c^P\) отображается в \(h\), поэтому
\(\langle c^P\rangle\le V\). Кроме того, \(c^V=c^N\):
каждый элемент \(V=\langle c\rangle N\) записывается \(c^k n\),
и \(c^{c^k n}=c^n\).
Нормальность \(N\) обеспечивает оба порядка разложения \(V\).

Для \(A_0=P\), \(B_0=V\) и
\[
 A_{j+1}=\langle c^{A_j}\rangle,\qquad
 B_{j+1}=\langle c^{B_j}\rangle
\]
монотонность даёт индукцией
\[
 A_{j+1}\le B_j\le A_j\qquad(j\ge0).
\]
База есть \(A_1\le V\le P\).
Применение монотонного оператора к двум включениям
даёт следующие два включения.
Если \(A_m=A_{m+1}=T\), то и \(A_{m+2}=T\);
сжатие даёт \(B_m=B_{m+1}=T\).
Теперь применяем лемму~\ref{lem:towers}.
\end{proof}

\begin{proposition}\label{prop:groupquandle}
Пусть \(\sigma:G\twoheadrightarrow H\), \(h=\sigma(c)\) и
\(H=\langle h^H\rangle\).
Если полный слой \(c^P\), \(P=\sigma^{-1}(C_H(h))\),
принадлежит \(\OS\), то \(c^G\in\OS\).
\end{proposition}
\begin{proof}
Пусть \(H'\) --- коммутант \(H\), \(N=\ker\sigma\),
\(L=\sigma^{-1}(H')\).
Из нормального порождения \(H\) элементом \(h\) следует
\[
 H=H'\langle h\rangle,\qquad G=L\langle c\rangle.
\]
Ограничение \(\sigma|_L\) сюръективно на \(H'\)
и \(L/N\cong H'\).
Положим \(\phi(l)=c^{-1}lc\), \(\psi(q)=h^{-1}qh\).
Так как \(N=\ker\sigma\), имеем \(N\trianglelefteq L\)
и \(\phi(N)=N\).
Индуцированный автоморфизм \(\bar\phi\) соответствует \(\psi\).

Подгруппа \(J=[H',\psi]\) нормальна в \(H'\)
и инвариантна относительно \(h\), поэтому нормальна в \(H\).
В \(H/J\) образ \(h\) централен и нормально порождает весь фактор.
Фактор цикличен; следовательно, \(H'\le J\), то есть
\begin{equation}\label{eq:headperfect}
 [H',\psi]=H'.
\end{equation}

Введём вспомогательные группы с бесконечными циклическими
генераторами \(u,v\):
\[
 \widetilde G=L\rtimes_\phi\langle u\rangle,\qquad
 \widetilde H=H'\rtimes_\psi\langle v\rangle,
 \qquad u^{-1}lu=\phi(l),\quad v^{-1}qv=\psi(q).
\]
Корректны эпиморфизмы
\[
 \widetilde\sigma(lu^k)=\sigma(l)v^k,\qquad
 \rho_G(lu^k)=lc^k,\qquad \rho_H(qv^k)=qh^k.
\]
Первое отображение имеет ядро \(N\), а остальные согласованы
с \(\sigma\) и \(\widetilde\sigma\).
Ограничения \(\rho_G,\rho_H\) дают изоморфизмы квандлов
\begin{equation}\label{eq:liftclasses}
 u^{\widetilde G}\cong c^G,\qquad
 v^{\widetilde H}\cong h^H.
\end{equation}
Действительно, сюръективность отображений групп позволяет
поднять любой сопрягающий элемент. Каждый элемент
\(u^{\widetilde G}\) имеет \(\Z\)-координату \(1\),
а на \(Lu\) отображение \(\rho_G\) инъективно:
\(lc=l'c\) влечёт \(l=l'\).
Для основания аргумент тот же.
Изоморфизмы \eqref{eq:liftclasses} совместимы с эпиморфизмами
квандлов и сохраняют все их полные слои;
также \(u^N\cong c^N\).

Для \(\phi\)-инвариантной \(X\le L\) нормальное замыкание \(u\)
в \(X\rtimes_\phi\langle u\rangle\) равно
\begin{equation}\label{eq:closureu}
 [X,\phi]\rtimes\langle u\rangle.
\end{equation}
Оно содержит \(u\) и все
\(x^{-1}\phi(x)=(u^x)^{-1}u\).
Обратно, подгруппа справа нормальна:
\([X,\phi]\trianglelefteq X\), она инвариантна относительно \(u\),
а на \(X/[X,\phi]\) действие \(u\) тождественно.
Следовательно, она содержит нормальное замыкание \(u\).

По лемме~\ref{lem:kernelfiber} \(c^N\in\OS\), значит,
\(u^N\in\OS\). При этом
\(u^{N\rtimes\langle u\rangle}=u^N\).
По \eqref{eq:closureu} ряд нормальных замыканий \(u\)
в этой группе имеет члены \(R_i\rtimes\langle u\rangle\),
где \(R_0=N\), \(R_{i+1}=[R_i,\phi]\).
Лемма~\ref{lem:towers} даёт стабилизацию \(R_i\).
Условие \eqref{eq:headperfect} позволяет применить
теорему~\ref{thm:group}; глобальный ряд
\(D_0=L\), \(D_{i+1}=[D_i,\phi]\) стабилизируется.
В силу \eqref{eq:closureu} стабилизируется и ряд нормальных
замыканий \(u\) в \(\widetilde G\).
Значит, \(u^{\widetilde G}\in\OS\), и
\eqref{eq:liftclasses} даёт \(c^G\in\OS\).
\end{proof}

\begin{remark}
Вспомогательное расщепление не является предположением
о расщеплении исходного \(\sigma\).
Элемент \(c\) может иметь конечный порядок;
генератор \(u\) выбран бесконечного порядка только во
вспомогательной группе.
\end{remark}

\section{Доказательство замкнутости для произвольных квандлов}
\label{sec:quandles}

\begin{lemma}\label{lem:connectedbase}
Пусть \(\pi:E\twoheadrightarrow C\) --- эпиморфизм,
\(C\) связен и все полные слои \(F_a=\pi^{-1}(a)\)
принадлежат \(\OS\). Тогда \(E\in\OS\).
\end{lemma}
\begin{proof}
Выберем произвольную первую орбиту \(O=\Orb(x,E)\).
По лемме~\ref{lem:images} \(\pi(O)=C\), а по
лемме~\ref{lem:slices} \(O\cap F_a\in\OS\) для всех \(a\in C\).

Положим \(G=\Inn(E)\), \(H=\Inn(C)\),
\(c=L_x^E\), \(h=L_{\pi(x)}^C\).
Эпиморфизм \(\pi\) индуцирует \(\sigma:G\twoheadrightarrow H\).
Отображение \(\lambda_E(y)=L_y^E\) является гомоморфизмом
квандлов, поскольку
\[
 L_{y\conj z}^E=L_y^E L_z^E(L_y^E)^{-1}.
\]
Эквивариантность этого отображения даёт
\(\lambda_E(O)=c^G\).
Связность \(C\) аналогично даёт
\[
 \{L_a^C:a\in C\}=h^H,\qquad H=\langle h^H\rangle.
\]

Обозначим \(b=\pi(x)\) и
\[
 T=\{a\in C:L_a^C=L_b^C\},\qquad
 S=O\cap\pi^{-1}(T).
\]
Квандл \(T\) тривиален: для \(a,d\in T\)
\(a\conj d=L_a^C(d)=L_d^C(d)=d\).
Отображение \(S\twoheadrightarrow T\) имеет полные слои
\(O\cap F_a\in\OS\); по лемме~\ref{lem:trivialbase}
\(S\in\OS\).

Полный слой \(c^G\to h^H\) над \(h\) точно равен
\(\lambda_E(S)\). В самом деле, образ элемента
\(L_y^E\), \(y\in O\), равен \(h\) тогда и только тогда,
когда \(L_{\pi(y)}^C=L_b^C\).
По лемме~\ref{lem:images} этот слой принадлежит \(\OS\).
Предложение~\ref{prop:groupquandle} даёт \(c^G\in\OS\).

Эпиморфизм \(\lambda_E|_O:O\twoheadrightarrow c^G\)
является накрытием: равенство глобальных трансляций
\(L_y^E=L_z^E\) влечёт равенство их ограничений
\(L_y^O=L_z^O\).
По лемме~\ref{lem:covering} \(O\in\OS\).
Это верно для каждой первой орбиты в \(E\);
следовательно, любой орбитальный ряд в \(E\) стабилизируется.
\end{proof}

\begin{proof}[Доказательство теоремы~\ref{thm:main}]
Возьмём произвольный ряд \(Q_i\) в \(Q\).
Образы \(B_i=\pi(Q_i)\) образуют орбитальный ряд в \(B\)
и после некоторого \(n\) равны связному под-квандлу \(B_*\).
Ограничение \(Q_n\twoheadrightarrow B_*\) имеет полные слои
\(Q_n\cap F_b\in\OS\) по лемме~\ref{lem:slices}.
Применяя лемму~\ref{lem:connectedbase}, получаем \(Q_n\in\OS\).
Поэтому выбранный хвост ряда стабилизируется.
Поскольку исходный ряд был произвольным, \(Q\in\OS\).
\end{proof}

\begin{remark}
Отображение \(y\mapsto L_y\) не обязано быть инъективным.
Таким образом, доказательство охватывает и nonfaithful-квандлы.
Общая конечная граница глубин полных слоёв нигде не требуется.
Граница в теореме~\ref{thm:group} относится к групповым рядам;
оптимальная количественная граница для произвольных
квандлов в этой работе не устанавливается.
\end{remark}

\section{\texorpdfstring{Отрицательный ответ для \(\OS_\omega\)}
{Отрицательный ответ для OSω}}\label{sec:negative}

\begin{proof}[Доказательство предложения~\ref{prop:negative}]
Пусть \(D_{2^r}\) --- диэдральный квандл на \(\Z/2^r\Z\)
с операцией \(a\conj b=2a-b\).
Рассмотрим дизъюнктную сумму
\[
 Q=\bigsqcup_{r\ge1}D_{2^r},
\]
где между разными слагаемыми \(x\conj y=y\).
Это квандл из \cite[пример~4.2]{BCNW}.
Проекция \(\pi:Q\to T_{\mathbb N_{\ge1}}\)
на номер слагаемого является эпиморфизмом на тривиальный квандл.
Её полные слои --- конечные \(D_{2^r}\), поэтому все они
принадлежат \(\OS_\omega\); основание принадлежит \(\OS_1\).

Глубины источника не ограничены.
Выберем нуль в слагаемом \(D_{2^r}\) и постоянный якорь \(0\).
Внутренние орбиты диэдрального квандла чётного порядка
суть классы чётности: \(L_aL_b\) есть трансляция на \(2(a-b)\).
Элементы остальных слагаемых действуют тождественно.
Поэтому ряд в \(Q\) имеет вид
\[
 Q_0=Q,\qquad
 Q_i=2^i\Z/2^r\Z\quad(1\le i\le r),\qquad
 Q_r=Q_{r+1}=\{0\}.
\]
Каждый шаг до \(r\) строгий.
Для любой предполагаемой общей границы \(n\) можно взять \(r>n\),
откуда \(Q\notin\OS_n\).
Следовательно, \(Q\notin\OS_\omega\).

Для любого косета \(S=a+2^j\Z/2^r\Z\) и \(x\in S\)
имеем
\[
 \Orb(x,S)=x+2^{j+1}\Z/2^r\Z.
\]
Действительно, \(L_y(x)=2y-x\) при \(y\in S\)
пробегает весь указанный косет, который инвариантен
относительно всех \(L_y^{\pm1}\).
Поэтому всякий отдельный ряд после первого шага находится
в одном слагаемом и, при любом выборе последующих якорей,
достигает одноэлементного квандла не позднее шага \(r\).
Значит, \(Q\in\tOS\subseteq\OS\), что согласуется
с теоремой~\ref{thm:main}.
\end{proof}

\begin{remark}\label{rmk:uniform}
Каждый слой последнего примера имеет конечную глубину, но общего
индекса для всех слоёв нет. Если определить иной вариант
замкнутости с общей границей глубин, этот пример его не опровергает.
Тот же пример показывает незамкнутость
\(\tOS_\omega=\bigcup_{n\ge1}\tOS_n\):
слои принадлежат \(\tOS_\omega\), а источник --- нет.
Следовательно, утверждение о замкнутости \(\tOS_\omega\)
в следствии~5.6 версии arXiv v3 работы~\cite{BCNW}
неверно для определения расширения из её раздела~2.2,
в котором единая граница по слоям не требуется.
Лемма~5.5 той же версии предполагает один фиксированный
индекс \(m\) для всех полных слоёв и этим примером
не опровергается. Для вывода о классе \(\tOS_\omega\)
из указанной леммы необходимо сохранить эту гипотезу.
\end{remark}

\begin{thebibliography}{9}
\bibitem{BCNW}
M.~Bonatto, A.~S.~Crans, T.~Nasybullov, G.~T.~Whitney.
\newblock Quandles with orbit series conditions.
\newblock \emph{Journal of Algebra}, \textbf{567} (2021), 284--309.
\newblock DOI:
\href{https://doi.org/10.1016/j.jalgebra.2020.09.026}
{10.1016/j.jalgebra.2020.09.026}.
\newblock Нумерация в тексте соответствует
\href{https://arxiv.org/abs/2004.10227v3}{arXiv:2004.10227v3},
18 августа 2020 года.
\end{thebibliography}

\end{document}
